\begin{lemma}
Let $V$ be finite and let $\mu,\nu$ be measures on
$\mathcal P(V)\times\mathbb R^V$ with discrete activation and product Borel
priority structures. Suppose $m_A<\infty$ for every $A\subseteq V$, and
both measures assign mass $m_A$ to $H_A=\{(B,\pi):B=A\}$.
Suppose also that, for $\rho=\mu,\nu$ and every $t:V\to[0,1]$,
\[
\rho\{(B,\pi):B=A,\ 0\le\pi_v\le t_v\ \forall v\}
=\rho(H_A)\operatorname{ofReal}(\prod_v t_v).
\]
Then $\mu=\nu$.
\end{lemma}
\begin{proof}
Write $Q=[0,1]^V$. For either measure $\rho$ and each $A$, setting all
thresholds equal to one gives $\rho(H_A\cap(\mathcal P(V)\times Q))=m_A$.
Both this set and $H_A$ are measurable, and $\rho(H_A)=m_A<\infty$.
Subtracting their finite masses gives
$\rho(H_A\setminus(\mathcal P(V)\times Q))=0$.
The finitely many $H_A$ partition the whole space. Hence both measures are
finite, with total mass $\sum_A m_A$, and both vanish outside
$\mathcal P(V)\times Q$.

Fix a measurable event $E$. For each $A$, let
$C_A=\{q:(A,q)\in E\}\cap Q$. This is Borel, because
$q\mapsto(A,q)$ is measurable and $Q$ is a finite intersection of closed
coordinate conditions. Apply
\noderef{SamplingFiniteProductLowerOrthantEqualityMeasurable} separately
to $\rho=\mu$ and $\rho=\nu$, with $F=H_A$, priority map the second
projection, $M=m_A$, and $C=C_A$. The projection is measurable; the atom
has finite mass $m_A$; and the assumed rectangle formula is exactly that
lemma's lower-rectangle hypothesis. Its conclusion, since $C_A\subseteq Q$,
is
\[
\rho\{(B,q):B=A,\ q\in C_A\}=m_A\operatorname{vol}(C_A).
\]
Thus the two measures agree on $E\cap H_A\cap(\mathcal P(V)\times Q)$.
The omitted part of $E\cap H_A$ is null by the first paragraph, so they agree
on $E\cap H_A$. Summing over the finite disjoint activation partition gives
$\mu(E)=\nu(E)$. Since this holds for every measurable $E$, the measures
are equal. If $V$ is empty, there is one activation atom and the same argument
uses the zero-dimensional cube and empty products, so no exception is required.
\end{proof}
